% Notes of the Montreal seminar, 1970 (cote 129, typed part): archive cote 129, site page 46 (PDF page 47).
% typescript. Draft compared: none.
% AI-transcribed by Claude Opus 5.5 (claude-opus-5-5, adaptive thinking, effort high), Message Batches API, batch pilotnext (msgbatch_01A9sdu8SYq23bEtqX9tUncQ), 2026-11-13; not yet checked by eye.
\scanpage{129/46}
\mspage{II.18}

Le foncteur $D^*$ recherché est alors donné par la composition
de toutes ses équivalences. La situation est résumée par le diagramme
suivant :

\diagram{
\begin{tikzcd}[column sep=large, row sep=large]
\left(\begin{array}{l}\text{p-groupes finis de}\\ \text{type multiplicatif}\end{array}\right)^{0}
\arrow[r, dashed, "D^*", "\hand{$\approx$}"']
\arrow[d, "\hand{$\wr$}"', "\text{dualité de Cartier}"]
&
\left(\begin{array}{l}\text{sous catégorie pleine de } \underline{\mathrm{Mod}}_{\hand{$\mathcal{D}$};\ W\ \hand{$\ell.f.$}}\\ \text{formée des objets où } V \text{ est } \deleted{objectif}\ \inserted{bijectif}\end{array}\right)
\arrow[d, "\hand{$\wr$}"]
\\
\left(\begin{array}{l}\text{p-groupes finis}\\ \text{étales}\end{array}\right)
\arrow[d, "\hand{$\wr$}"', "\text{dualité de Pontryagin}"]
&
\left(\begin{array}{l}\text{W-modules de longueur finie munis d'un}\\ \text{isomorphisme } V : M \to M^{\sigma}\end{array}\right)
\arrow[d, "\hand{$\wr$}"]
\\
\left(\begin{array}{l}\text{p-groupes finis}\\ \text{étales}\end{array}\right)^{0}
\arrow[r, "D^*", "\hand{$\sim$}"']
&
\left(\begin{array}{l}\text{sous catégorie pleine de } \underline{\mathrm{Mod}}_{\hand{$\mathcal{D}$};\ W\ \hand{$\ell.f.$}}\\ \text{formée des objets où } F \text{ est bijectif}\end{array}\right)
\end{tikzcd}
}
\margin{X}

(4.3.3)

4.5. \uline{Cas général}

Si $G$ est un p-groupe fini quelconque, on écrit,(d'après 3.2)
\[ G = G^{\mathrm{uni}} \times G^{\mathrm{mult}} \]
où $G^{\mathrm{uni}}$ et $G^{\mathrm{mult}}$ sont respectivement unipotent et de type multiplica-%
tif, et on pose
\[ D^*(G) = D^*(G^{\mathrm{uni}}) \times D^*(G^{\mathrm{mult}}) \ . \]

Le théorème de Dieudonné 4.2. sera démontré si l'on prouve que
l'on a une équivalence de catégories

\mspage{46}
